\documentclass[conference]{IEEEtran}
\usepackage{cite}
\usepackage{amsmath,amssymb,amsfonts}
\usepackage{graphicx}

\begin{document}

\title{An Observation-Driven Digital Twin for Hybrid Smart Parking Operations}

\author{\IEEEauthorblockN{1\textsuperscript{st} Ashutosh Sononey}
\IEEEauthorblockA{\textit{Computer Engineering} \\
\textit{Fr. C. Rodrigues Institute of Technology}\\
Mumbai, India \\
ashutoshsononey@gmail.com}}

\maketitle

\begin{abstract}
This paper presents the digital-twin subsystem of PragmaPark, a hybrid smart-parking platform that combines Internet of Things sensing, occupancy prediction, reinforcement-learning pricing, blockchain audit logging, and physical actuation. The subsystem maintains a durable record of lot observations, derived states, forecasts, scenario runs, and model provenance. Its operational forecast is a persistence baseline anchored to a timestamped observation and evaluated only against a later observation within a horizon-specific tolerance. The platform also provides deterministic parking scenarios with empirical variability bands derived from historical occupancy transitions. These scenarios are decision-support tools rather than causal estimates and cannot change live prices or actuators. An experimental spatial-temporal forecaster and a conditional generative model are retained for offline investigation but are not deployed as production forecasting or intervention models. The resulting design provides an auditable measurement loop while preserving a clear boundary between observed evidence, simulation, and control.
\end{abstract}

\begin{IEEEkeywords}
digital twin, smart parking, occupancy forecasting, counterfactual analysis, dynamic pricing, intelligent transportation systems
\end{IEEEkeywords}

\section{Introduction}
Parking operations combine a limited physical resource with uncertain and time-varying demand. An operator must decide whether to alter price, reserve capacity, or redirect drivers before the consequences of that decision are directly observable. A digital twin can support these decisions by maintaining a virtual representation of an operational system and relating forecasts or simulations to subsequent observations. Existing smart-parking twin work ranges from sensor-backed representations of car parks to building-information-modeling and city-scale mobility systems \cite{b2,b3,b13}. Recent research has also considered spatial-temporal forecasting and generative methods for parking and urban-mobility analysis \cite{b1,b20}.

PragmaPark is designed as a hybrid platform rather than a digital twin in isolation. Its six layers are Internet of Things (IoT) sensing, supervised machine-learning prediction, blockchain logging, reinforcement learning (RL), digital twinning, and actuation. The twin has a distinct role within this pipeline. The supervised prediction service estimates occupancy from its engineered feature set, and the RL service selects a pricing action. The twin records operational evidence, creates traceable forecasts, and provides a separate environment for examining the consequences of explicit assumptions. It does not replace the pricing controller, and a scenario result is never an actuator command.

This paper describes the implemented subsystem and its operational boundaries. The central contribution is not a claim of city-scale predictive fidelity. Instead, it is a durable and auditable architecture in which forecast inputs, outputs, outcomes, scenario assumptions, and model versions can be inspected after the fact.

\section{System Position in the Hybrid Platform}

The IoT layer fuses parking evidence from the platform's sensor interfaces. The machine-learning layer produces occupancy and price-related predictions from historical temporal and parking-event features. Blockchain services provide a tamper-evident audit ledger for operational transactions. The RL layer selects bounded price adjustments, while actuators apply approved operational changes such as barrier, pricing-board, and congestion-light updates. Dynamic pricing is naturally modelled as a sequential decision problem because a price action affects both utilization and revenue \cite{b11}. Multi-agent approaches have similarly been investigated for parking assignment across multiple areas \cite{b23}.

The digital twin intersects these layers without collapsing their responsibilities. A completed parking session provides an operational occupancy and price observation. The twin persists this observation and creates baseline forecasts. A manager may then request a scenario, but that request creates only a recommendation record. The RL policy and the actuator bridge remain the paths through which a production action is selected and applied.

\subsection{Execution Flow}

The persistent twin has two current evidence-entry paths. The session-completion path writes an observation with source \texttt{session\_end}, derives a state record, and creates forecasts for the configured horizons. The administrative application programming interface accepts the same canonical observation shape for controlled ingestion or historical replay. Both paths lead to the same persistence, forecast, evaluation, and metrics services. Scenario requests follow a separate path: a base state and named assumption are passed to the deterministic engine, then the result is written only as a \texttt{TwinScenarioRun}. It has no write path to a parking lot, reservation, price board, barrier, or congestion light.

The raw IoT ingestion route is not yet automatically subscribed to the persistent twin service. Consequently, direct sensor traffic becomes twin evidence only when it is propagated through an authorized canonical observation request or an operational session-completion event. This boundary is an implementation limitation, not a claim that all sensor readings are already incorporated into the twin.

\section{Persistent Twin Data Model}

Earlier prototypes represented twin state in process memory. The current implementation uses database-backed records so that evidence and forecasts survive a service restart. Five entities capture the operational history.

\begin{itemize}
\item \texttt{TwinObservation} stores a timestamped lot observation, including occupied and total slots, occupancy rate, price, arrivals, departures, source, confidence, and contextual metadata.
\item \texttt{TwinState} records the state derived from an observation, including estimated occupancy, availability, price, confidence, and the source observation identifier.
\item \texttt{TwinForecast} stores an immutable forecast, its source observation, generation time, target time, horizon, model name, model version, interval bounds, and later evaluation fields.
\item \texttt{TwinScenarioRun} stores a read-only what-if result with its assumptions, uncertainty note, random seed when applicable, and execution latency.
\item \texttt{TwinModelVersion} records an artifact version, feature schema, training-data cutoff, validation metadata, and promotion status.
\end{itemize}

All timestamps are normalized to UTC before persistence. The observation timestamp, rather than the wall-clock time of an API request, anchors a forecast. This avoids assigning historical knowledge to a forecast that could not have used it.

For lot $l$ at time $t$, an observation is represented by
\begin{equation}
x_{l,t}=(o_{l,t},n_{l,t},C_l,p_{l,t},a_{l,t},d_{l,t},q_{l,t},c_{l,t}),
\label{eq:observation}
\end{equation}
where $o$ is occupancy rate, $n$ is occupied slots, $C$ is capacity, $p$ is the current price, $a$ and $d$ are arrivals and departures, $q$ is sensor confidence, and $c$ is contextual metadata. The service derives availability as $C_l-n_{l,t}$ and persists a state record linked to the observation. This representation separates a measured input from a forecast or a simulated scenario result.

\section{Forecast Lifecycle}

The service always produces a persistence forecast. For a source observation with occupancy $o_t$ and horizon $h$, the baseline is
\begin{equation}
\hat{o}_{t+h}=o_t.
\label{eq:persistence}
\end{equation}

This baseline is intentionally simple. It provides a reference that a more complex model must beat on later observed outcomes before any stronger claim is justified. A generated forecast is linked to its source observation and is not overwritten when its outcome becomes available. Instead, the later observed occupancy, signed error, absolute error, and evaluation timestamp are attached to the original record.

The evaluator selects the first observation at or after the target time only when it falls within a bounded tolerance for that horizon. This prevents a 15-minute forecast, for example, from being scored against a reading acquired many hours later. Metrics are computed separately for each lot, model version, and horizon. The reported quantities include mean absolute error, root mean squared error, signed bias, interval coverage, calibration error, skill relative to persistence, evaluated-forecast fraction, and a simple recent-versus-earlier error drift measure.

For $N$ evaluated forecasts at a fixed lot and horizon, the primary error measures are
\begin{equation}
\operatorname{MAE}=\frac{1}{N}\sum_{i=1}^{N}|o_i-\hat{o}_i|,
\qquad
\operatorname{RMSE}=\sqrt{\frac{1}{N}\sum_{i=1}^{N}(o_i-\hat{o}_i)^2}.
\label{eq:forecastmetrics}
\end{equation}
The persistence baseline makes the skill score for a candidate model $m$ explicit:
\begin{equation}
\operatorname{Skill}_m=1-
\frac{\operatorname{MAE}_m}{\operatorname{MAE}_{\mathrm{persistence}}}.
\label{eq:skill}
\end{equation}
A positive value indicates improvement over persistence on the same evaluated population; it does not by itself establish production readiness.

The persistence service currently stores a fixed interval of $\hat{o}_{t+h}\pm0.05$, clipped to $[0,1]$. The service reports empirical coverage of this interval and the absolute difference between that coverage and 0.90. These quantities are monitoring statistics; the fixed interval is not a statistically calibrated predictive interval.

Completed sessions are one live source for this lifecycle. When a session ends, PragmaPark records the resulting occupancy and rate as a \texttt{TwinObservation}, derives a \texttt{TwinState}, and creates baseline forecast rows. The administrative twin application programming interface (API) can also accept a canonical observation record. These endpoints require administrator authorization because operational observations must not be publicly writable.

\section{Scenario Analysis}

The scenario engine contains six explicit rule-based projections: zone closure, price surge, capacity expansion, weather disruption, holiday demand, and resident-share adoption. Table~\ref{tab:scenarios} summarizes their current semantics. The constants are transparent policy assumptions, not fitted intervention effects.

\begin{table}[!t]
\caption{Implemented Parking Scenarios}
\label{tab:scenarios}
\centering
\footnotesize
\begin{tabular}{p{0.26\columnwidth}p{0.64\columnwidth}}
\hline
\textbf{Scenario} & \textbf{Projection} \\
\hline
Zone closure & Removes availability, sets occupancy to full, and applies a 1.5 price multiplier. \\
Price surge & Applies a 1.5 price multiplier and subtracts 0.15 from occupancy rate. \\
Capacity expansion & Increases capacity by 20\% and multiplies occupancy rate by 0.83. \\
Weather disruption & Applies a fixed 0.30 reduction in occupancy rate. \\
Holiday spike & Multiplies occupancy by 1.25 and price by 1.10. \\
Resident-share adoption & Releases the rounded value of 15\% of total slots from occupied slots under an explicit sharing assumption. \\
\hline
\end{tabular}
\end{table}

The calibrated scenario endpoint wraps a deterministic projection with an empirical variability interval. For the requested lot and horizon, it samples individual historical occupancy changes and adds them to the scenario point estimate. Consequently, the interval describes the range of ordinary observed variation around an assumed scenario result. It does not estimate the causal effect of a price increase, a closure, or weather. Runs supported by fewer than 20 historical transitions are marked experimental; with no matching history, the component returns the broad bounded interval $[0,1]$ rather than inventing evidence.

More precisely, if $g_s(x_{l,t})$ is the deterministic result for scenario $s$ and $\Delta_{l,h}$ is the empirical distribution of observed occupancy changes for lot $l$ and horizon $h$, the service reports
\begin{equation}
\hat{o}^{(s)}_{l,t+h}=\operatorname{clip}(g_s(x_{l,t}),0,1),
\label{eq:scenario}
\end{equation}
and an empirical interval
\begin{equation}
\left[\operatorname{clip}(\hat{o}^{(s)}_{l,t+h}+Q_{0.05}(\Delta_{l,h}),0,1),
\operatorname{clip}(\hat{o}^{(s)}_{l,t+h}+Q_{0.95}(\Delta_{l,h}),0,1)\right].
\label{eq:scenariointerval}
\end{equation}
The quantiles characterize observed variability, not a scenario-specific treatment distribution.

This distinction is important for evaluation. A later observation from an ordinary day cannot validate a hypothetical closure or price surge. The system therefore records scenario latency but reports scenario backtest error as unavailable until a dataset contains labelled interventions and their later outcomes. This avoids converting unrelated future observations into spurious causal evidence.

\section{Experimental Forecasting Components}

PragmaPark retains two experimental components outside the live twin path. The first is a compact spatial-temporal identity predictor. It uses hour-of-day and day-of-week embeddings together with a real-distance adjacency signal constructed from lot coordinates and, where available, road-network distance. Its training helper reads timestamped \texttt{OccupancyRecord} history chronologically, using a later record as the target. That legacy record type does not carry the source-provenance field of \texttt{TwinObservation}; therefore, the helper assumes its input history represents operational observations and should not be used as evidence of source-validated sensor learning. It is not automatically retrained or automatically promoted by the live service. Promotion requires comparison with the persistence baseline on held-out evaluated outcomes.

The second component is a conditional variational autoencoder with a Wasserstein generative adversarial critic. Its training representation includes normalized occupancy, price, congestion, duration, and resident-share information. It is retained as a versioned offline research artifact. It is not loaded by the persistent twin service and does not drive live forecasts, scenario results, prices, or actuators. In particular, the project has no labelled intervention-to-outcome dataset from which it could learn causal effects for closure, pricing, capacity, weather, holiday, or sharing interventions. It must therefore not be interpreted as a deployed causal model.

The legacy in-memory simulator remains available for isolated what-if exploration. It uses a price-response rule and noise to advance a simulated state, but its output is not written as an observation and is not used as a training target for the persistent forecast path.

\section{Operational Boundaries and Limitations}

The implemented twin is measurement-ready, not yet a validated city-scale forecasting system. The persistence baseline is operational and auditable, but a stable published error bound requires a sufficiently large stream of representative observed outcomes. The experimental spatial-temporal model has no automatic training schedule or production promotion path. Its spatial input is based on lot proximity rather than a calibrated model of driver substitution or traffic flow, and its current training helper requires stronger observation provenance before it can support a source-validated learning claim.

Similarly, the deterministic scenarios are useful because their assumptions are inspectable, but they are not estimates of causal treatment effects. Their empirical variability bands quantify historical occupancy movement, not the effect of the proposed intervention. Labelled before-and-after intervention records are required before scenario error or causal performance can be measured honestly.

The architecture also separates simulation from control. Twin scenarios cannot change a lot state, price, reservation, or actuator. Any production action must be selected through the appropriate policy or authorized operational workflow. This separation reduces the risk that an exploratory output is mistaken for an approved command.

\section{Validation Protocol and Current Evidence}

The implementation supports validation at two different levels. First, automated tests verify the architecture: an observation produces a persisted state; forecasts retain their source observation and model version; only horizon-aligned later observations may populate evaluation fields; forecasts remain immutable after evaluation; process restart preserves the database records; scenario execution does not alter production state; and authorization protects twin endpoints. These tests establish lifecycle and safety properties, not predictive accuracy.

Second, the persistent forecast records support empirical validation as operational observations accumulate. A valid evaluation requires a forecast made before its target time and a later matched observation within the service tolerance. The system can then aggregate the measures in \eqref{eq:forecastmetrics} by lot, model, and horizon, compare candidate models through \eqref{eq:skill}, and inspect interval coverage and drift over time. The project does not presently report a stable field accuracy bound because its retained operational history is insufficient for a representative estimate.

Scenario validation has a stricter requirement. It needs records that identify an intervention, its scope and timing, the pre-intervention state, and a later observed outcome. Such labelled intervention-to-outcome data are not yet available in the project. Therefore, deterministic scenario projections and their intervals are validated for persistence, reproducibility, and non-actuation, but are not reported as empirically validated causal predictions.

\section{Conclusion}

This paper has described an observation-driven digital twin for a hybrid smart-parking platform. The subsystem persists operational observations, derived state, immutable forecasts, scenario runs, and model provenance. Its live forecasting path begins with an explicit persistence baseline whose outputs can be evaluated against horizon-aligned future observations. Its scenario engine provides transparent rule-based projections with empirical variability intervals while avoiding unsupported causal claims. Experimental spatial-temporal and generative components remain offline until they can be evaluated on adequate observed data. By keeping evidence, forecasting, simulation, and control separate, the design provides a practical foundation for measurable smart-parking decision support.

\section{Future Work}

The next step for the twin is not to add an untested model, but to connect and evaluate the models already present in PragmaPark. The platform's supervised occupancy pipeline contains Random Forest, XGBoost, and Ridge meta-model artifacts trained on parking history. A future twin adapter can register these artifacts as candidate \texttt{TwinForecast} producers, using the same source observation, target horizon, and immutable evaluation lifecycle as the persistence baseline. They should be compared by lot and horizon against persistence before any model is selected for operational use.

The STID candidate requires a more complete operational path: raw IoT ingestion should create provenance-bearing \texttt{TwinObservation} records; scheduled training should consume only those records; fitted artifacts should be versioned and reloadable; and promotion should require an adequate chronological held-out evaluation. This would turn the current experimental code into a measurable spatial-temporal forecasting candidate rather than an offline helper over legacy occupancy history.

The CVAE-WGAN requires an even higher evidential threshold. It should remain offline until intervention records contain intervention type, affected lot or slots, start and end time, pre-intervention state, and subsequent observed outcome. Such data would permit evaluation of a conditional model of future state given the current state and an intervention. Without those labels, the deterministic scenario engine remains the appropriate operational component, and the generative artifact should not be connected to pricing, actuation, or live scenario claims.

Finally, the twin should receive direct canonical observations from authenticated IoT ingestion, retain exogenous weather and event features in a model-ready schema, and report forecast and interval metrics on a scheduled dashboard. These changes would make the existing hybrid architecture observable end to end while maintaining the separation between forecasting, simulation, and production control.


\begin{thebibliography}{00}
\bibitem{b2} Y. Zou, F. Ye, A. Li, M. Munir, E. Hjelseth, and S. F. Sujan, ``A digital twin prototype for smart parking management,'' in \textit{ECPPM 2022}. CRC Press, 2023.
\bibitem{b3} A. A. Amusan and G. A. Ogunleye, ``A digital twin-enabled smart car park management system: architecture and impact on emission reduction,'' \textit{FUOYE Journal of Engineering and Technology}, vol. 9, no. 4, pp. 629--635, 2024, doi: 10.4314/fuoyejet.v9i4.10.
\bibitem{b13} J. K. Coching, R. K. Billones, R. K. C. Billones, A. K. M. Brillantes, S. Yunus, V. A. Pitogo, and R. Senkerik, ``Digital twinning mechanism and building information modeling for a smart parking management system,'' \textit{Smart Cities}, vol. 8, no. 5, p. 146, 2025, doi: 10.3390/smartcities8050146.
\bibitem{b1} F. Piccialli, S. Amitrano, D. Cerciello, A. Borrelli, E. Prezioso, and M. Canzaniello, ``A digital twin framework for urban parking management and mobility forecasting,'' \textit{Nature Communications}, vol. 16, no. 1, p. 9400, 2025, doi: 10.1038/s41467-025-65306-w.
\bibitem{b20} M. Canzaniello, S. Amitrano, E. Prezioso, F. Giampaolo, S. Cuomo, and F. Piccialli, ``Leveraging digital twins and generative AI for effective urban mobility management,'' in \textit{2024 IEEE Cyber Science and Technology Congress}, 2024, pp. 146--153, doi: 10.1109/CyberSciTech64112.2024.00032.
\bibitem{b11} L. Z. Poh, T. Connie, T. S. Ong, and M. K. O. Goh, ``Deep reinforcement learning-based dynamic pricing for parking solutions,'' \textit{Algorithms}, vol. 16, no. 1, p. 32, 2023, doi: 10.3390/a16010032.
\bibitem{b23} X. Zhang, C. Zhao, F. Liao, X. Li, and Y. Du, ``Online parking assignment in an environment of partially connected vehicles: a multi-agent deep reinforcement learning approach,'' \textit{Transportation Research Part C: Emerging Technologies}, vol. 138, p. 103624, 2022, doi: 10.1016/j.trc.2022.103624.
\end{thebibliography}

\end{document}
